\section{Synchronization and Coordination Problems}
\label{sec:problems}

This section states the problems the project set out to solve. The mechanisms that answer
them are in \cref{sec:addressing}, and the ones that answer failures are in
\cref{sec:fault-tolerance}.

Six of them are numbered P1 to P6 and are referred to by those names throughout the
document.

\subsection{The two exclusions in the system}
\label{sec:two-exclusions}

Two resources in this system are exclusive:

\begin{table}[h]
\centering\small
\caption{The two exclusions.}
\label{tab:exclusions}
\begin{tabularx}{\textwidth}{@{}l >{\raggedright\arraybackslash}p{3.2cm}
                               >{\raggedright\arraybackslash}p{4.2cm} X@{}}
\toprule
& Object protected & Contended by & Scope \\
\midrule
P1 & the connector & drivers, among themselves & inside one station \\
P2 & the vehicle & stations, among themselves & across nodes \\
\bottomrule
\end{tabularx}
\end{table}

A connector belongs to exactly one station, so the station decides alone. The vehicle instead, could appear at two stations at the same time, since the driver could open two browser
tabs and presses \emph{reserve} on both, each station finds its own connector free and
grants, and neither can detect the breach, because the fact that would reveal it is held by
neither. The second exclusion is the reason this project has a coordination
problem at all.

\subsection{Problems inside one station}
\label{sec:local-problems}

\paragraph{P1: several drivers want the same connector.} Two requests for connector 3 arrive
in the same millisecond. Both read the connector as free, both proceed, and the outlet ends
up promised twice. This is the classic race on a shared resource, and its solution is
\cref{sec:p1}.

\paragraph{P5: the site cannot power everything at once.} The sum of what the connectors can
deliver exceeds what the site's grid connection can supply, which is true of real
installations and is the reason dynamic load management exists.

A connector
is exclusive and is therefore \emph{granted}. Electrical capacity is divisible and is
therefore \emph{apportioned}, with the apportionment changing every time a car arrives,
leaves, or slows down as its battery fills. See
\cref{sec:p5}.

\subsection{Problems between the client and the system}
\label{sec:client-problems}

\paragraph{P3: somebody stops answering.} A driver reserves a connector and never arrives,
or disappears in the middle of a session with the cable still in the socket. The same problem in a second form: a station holds a claim at the
coordinator and dies, and the coordinator has no way to ask it whether the claim still
matters. See \cref{sec:p3}.

\paragraph{P6: everyone must see the same connector.} Two drivers, the charge point and the
back office all display a connector whose state changes several times a minute. The station
pushes to all of them (\cref{sec:ssr}); what is left to decide is what it pushes. Sending
only the change is smaller, but it leaves a client that misses one message wrong from then
on, with nothing to tell it so. See \cref{sec:p6}.

\subsection{Problems across nodes}
\label{sec:distributed-problems}

\paragraph{P2: one vehicle, one reservation, network-wide.} The scenario of
\cref{sec:two-exclusions}: the same car reserving at two stations at once, with each station
locally correct. See \cref{sec:p2}.

\paragraph{P2b: the authority itself fails.} Once P2 is answered by an authority, that
authority becomes the thing that can fail. If it
crashes, nobody can reserve until something replaces it. If a partition splits the network,
both sides may conclude that the other is gone and each appoint its own authority, and two
authorities granting claims for the same vehicle break exactly the guarantee the authority
exists to provide. "Split brain" is worse than downtime here, because downtime refuses and
split brain grants. See \cref{sec:p2b}.

\paragraph{P4: a station node goes away.} It can crash, or its uplink can be cut.

A crashed node takes its sessions with it: the process metering them is gone and no row
reaches the database. A partitioned node keeps delivering power to the cars plugged into it
while the rest of the network stops hearing from it, so the site works and the operator goes
blind. In both cases the vehicles that station was holding are stranded: their claims are
alive in the coordinator, and their drivers cannot reserve anywhere else until those claims
expire. See \cref{sec:station-failures}.
